\documentclass[a4paper]{article}
\usepackage{amsfonts}
\usepackage{amsmath}
\usepackage{amssymb}
\usepackage{graphicx}     

\begin{document}
  
\noindent \large Auteur: Kyan Van den Eynde \\
\noindent \large Studierichting: Informatica\\
\noindent \large Oefeningenreeks 6A nr. 2\\

\medskip

\normalsize

In een rij met algemene term $x_n = n^2 + n$ laat men de eerste term weg.\\

Wat wordt de nieuwe formule voor de $n$-de term?\\

$x_{n+1} = (n+1)^2+(n+1) = n^2 + 2n + 1 + n + 1 = n^2 + 3n + 2$\\

$x_n = n^2 + 3n + 2$\\

\begin{thebibliography}{999}
%\bibitem{a} Referentie
\end{thebibliography}
\end{document} 
